In simulated attitude panels of SHP size, a fixed classifier on three pooled long-lag correlations distinguished the declared network-influence law from trait-plus-autoregressive persistence matched to the same short-term stability. It first reached an estimated AUC of \resultnum{v4.protocol.target}{0.70} at nominal innovation-variance fraction $a=\resultnum{j.b.r055.first.point}{0.10}$ for signal shares from \resultnum{v4.tab.T1.B.r055.r}{0.55} to \resultnum{v4.tab.T1.B.r090.r}{0.90}, and at $a=\resultnum{j.b.r050.first.point}{0.30}$ for $r=\resultnum{v4.tab.T1.B.r050.r}{0.50}$. This describes ranking panels from two simulated laws under known reliability and population-moment matching, without establishing causal attribution or a threshold for an empirical item. The confirmatory run used new seeds after partial exposure to calibration results, mainly at $r=\resultnum{v4.protocol.r.070}{0.70}$.

The exploratory analyses demonstrate why that discrimination is fragile: response-channel misspecification alone separated panels without influence, while a single fit that moved B's signal share by \resultnum{j.x3.free.r_shift.min}{0.0011} to \resultnum{j.x3.free.r_shift.max}{0.0028} below the true value and matched one more lag absorbed most separation on later lags. For comparison of scales, three-wave estimates for the SHP trust item have standard errors of \resultnum{j.shp.heise.w2004.se}{0.011}--\resultnum{j.shp.heise.w2002.se}{0.014}. Effects of richer persistence and network burn-in remain untested. The demonstrated contribution is conditional classifier discrimination, together with specific measurement and matching sensitivities that an empirical influence interpretation must address.
